\documentclass[12pt]{article}
\usepackage[utf8]{inputenc}
\usepackage[T1]{fontenc}
\usepackage[margin=1in]{geometry}
\usepackage{setspace}
\usepackage{xcolor}
\usepackage[colorlinks=true, urlcolor=blue!60!black, linkcolor=blue!60!black]{hyperref}
\setlength{\parindent}{0em}
\setlength{\parskip}{0.45em}
\renewcommand{\baselinestretch}{1.0}
\begin{document}
\begin{center}
{\Large \textit{Chae Chan Ping v. United States}, 130 U.S. 581 (1889)\par}
\vspace{0.3em}
{\normalsize Case Brief --- One-Page Edition\par}
\end{center}

\textbf{Memory Jogger.} \textit{The Chinese Exclusion Case}: the later-in-time rule and the birth of the plenary power over immigration (\href{https://supreme.justia.com/cases/federal/us/130/581/}{130 U.S. 581}).

\textbf{Facts.} A Chinese laborer who left San Francisco with a statutory reentry certificate was refused landing because the Scott Act of 1888, passed while he was at sea, voided all certificates despite treaty guarantees.

\textbf{Procedural History.} The federal circuit court denied habeas corpus.

\textbf{Votes.} Affirmed, 8--0: Justice Field for a unanimous Court.

\textbf{Holding.} A later statute abrogates a conflicting treaty as domestic law, and Congress's inherent power to exclude noncitizens is conclusive upon the Judiciary.

\textbf{Reasoning.} Statutes and treaties are equally supreme, so ``the last expression of the sovereign will must control.'' Exclusion is an incident of nationhood, not of any enumerated clause, and Field deferred to Congress's view that ``foreigners of a different race'' were dangerous. A reentry certificate was a revocable license.

\textbf{Examples.} \emph{Same-Side:} a statute voiding treaty-promised visas. \emph{Opposite-Side:} hard labor without a jury (\textit{Wong Wing}). \emph{Fence-Sitter:} a religious exclusion statute.

\textbf{Critique.} Inherent sovereignty escapes enumerated-powers limits, and Field's racial reasoning shows the peril of deference that critics see continuing in \textit{Trump v. Hawaii}.

\textbf{Key Quotation.} ``For local interests the several states of the Union exist, but for national purposes, embracing our relations with foreign nations, we are but one people, one nation, one power. ... If, therefore, the government of the United States, through its legislative department, considers the presence of foreigners of a different race in this country, who will not assimilate with us, to be dangerous to its peace and security, their exclusion is not to be stayed because at the time there are no actual hostilities with the nation of which the foreigners are subjects. ... In both cases its determination is conclusive upon the judiciary.'' --- Justice Field, Opinion of the Court, 130 U.S.\ at 606.

\end{document}
